\begin{textsample}{POS Dim 2 – rr – Score 35.00 – rr/rr\_em\_3.txt}  \label{ex:f2_pos_144}
de ocupação ou aperfeiçoamento posteriores ; III – o aprimoramento do educando como pessoa humana, incluindo a formação ética e o desenvolvimento da autonomia intelectual e do pensamento crítico ; IV – a compreensão dos fundamentos científico-tecnológicos dos processos produtivos, relacionando a teoria com a prática, no ensino de cada disciplina.
Nessa perspectiva, o DCRR visa promover a permanência e orientar as aprendizagens dos estudantes roraimenses, respondendo às suas demandas e aspirações.
Portanto, com o advento da Lei nº 13.415 ( BRASIL, 2017a ), ficou estabelecido uma nova organização curricular do Ensino Médio que valoriza a formação integral dos sujeitos, o protagonismo juvenil e a construção dos projetos de vida, a partir de uma estrutura \textbf{flexível}, conforme define o Art. 36 : > O \textbf{currículo} do Ensino Médio será composto pela Base Nacional Comum Curricular e por \textbf{itinerários} formativos, que deverão ser organizados por meio da \textbf{oferta} de diferentes arranjos curriculares, conforme a relevância para o contexto local e a possibilidade dos sistemas de ensino a saber : > > I - linguagens e suas tecnologias ; > II - matemática e suas tecnologias ; > III - ciências da natureza e suas tecnologias ; > IV - ciências humanas e sociais aplicadas ; > V - formação \textbf{técnica} e profissional.
Além da nova estrutura curricular, a Lei nº 13.415 ( BRASIL, 2017a ) também \textbf{instituiu} a ampliação da carga-horária de 800h para 1.000h anuais a partir de 2022, assim como, a Política de Fomento à Implementação de Escolas de Ensino Médio em Tempo Integral, regulamentada pela \textbf{Portaria} MEC nº 727 ( BRASIL, 2017b ) e \textbf{atualizada} pela \textbf{Portaria} MEC nº 2.116 ( BRASIL, 2019b ), com previsão de repasse de recursos pelo período de 10 anos.
Ainda na intenção de garantir a implementação das mudanças definidas em lei, o Ministério da Educação – MEC, em diálogo com o Conselho Nacional de \textbf{Secretários} de Educação – CONSED, elaborou o Programa de Apoio ao Novo Ensino Médio, com o objetivo de subsidiar as 27 unidades federativas na elaboração e execução de um Plano de Implementação do Novo Ensino Médio, contribuindo para atingir as metas do Plano Nacional de Educação – PNE 2014/2024 ( BRASIL, 2014 ).
Considerando estas inovações, é importante ressaltar que o Documento Curricular de Roraima – DCRR, para o Ensino Médio, foi elaborado em regime de colaboração entre estado e \textbf{municípios}, por meio do termo de Adesão entre o Conselho Nacional de \textbf{Secretários} de Educação - CONSED e União Nacional dos \textbf{Dirigentes} \textbf{Municipais} de Educação – UNDIME.
Desta forma, O DCRR tem a responsabilidade de representar o território em sua diversidade de saberes e vivências culturais, fortalecendo o desenvolvimento de projetos de vida, ao possibilitar escolhas de estilos de vida saudáveis, sustentáveis e éticas.
Ao pensar os princípios norteadores do Ensino Médio, é preciso discutir meios para o direcionar o esforço pedagógico para adoção de metodologias adequadas às necessidades e interesse dos alunos, organizadas nas práticas didáticas dos \textbf{currículos} escolares.
Estas mudanças ampliam as perspectivas de pensar a formação continuada dos profissionais da educação, para o desenvolvimento de metodologias didáticas e pedagógicas, fundamentadas nos processos, procedimentos inerentes à realidade local e ao papel docente na sala de aula da Educação Básica.
Nesse sentido, o estado de Roraima, por meio da Secretaria de Estado da Educação e Desporto – SEED, aderiu ao Programa de Apoio ao Novo Ensino Médio, \textbf{instituído} por meio da \textbf{Portaria} nº 649 ( BRASIL, 2018b ), que tem como objetivo dar suporte às unidades da federação na elaboração e execução do Plano de Implementação do Novo Ensino Médio, que contemple a Base Nacional Comum Curricular nas escolas de Ensino Médio do país.
O Programa de Apoio ao Novo Ensino Médio tem como objetivos : I. apoiar as \textbf{secretarias} de educação estaduais e do Distrito Federal para que adaptem seus \textbf{currículos}, contemplando a BNCC, \textbf{itinerários} formativos e ampliação da \textbf{carga} \textbf{horária} mínima para 3.000 horas ; II. fortalecer as escolas de Ensino Médio em Tempo Integral nos estados e Distrito Federal ; III. apoiar os estados e o Distrito Federal na melhoria do monitoramento e avaliação de suas \textbf{políticas} e programas, aprimorando sua capacidade de \textbf{gestão} ; IV. criar mecanismos de responsabilização e pactuação de resultados entre os entes federados, garantindo maior apoio às redes mais vulneráveis ; V. fomentar mecanismos de mobilização e compartilhamento de melhores práticas entre as redes e entre as escolas no sentido de otimizar a implementação do Programa ; VI. apoiar a implementação do Novo Ensino Médio, promovendo o acesso a \textbf{itinerários} formativos de forma equitativa, tanto da perspectiva socioeconômica, quanto geográfica, etnicorracial e de gênero.
Para efetivação desses objetivos, a Resolução CNE nº 4, ( BRASIL, 2018c ), estabelece que os \textbf{currículos} devem assegurar, de forma integrada, uma formação geral básica com seus \textbf{itinerários} formativos seguindo os termos das \textbf{Diretrizes} Curriculares Nacionais para o Ensino Médio - DCNEM ( BRASIL, 2018d ).
No \textbf{artigo} 10, a Formação Geral Básica é contemplada como referência nacional obrigatória tendo como base a ampliação de conhecimentos por meio do desenvolvimento e consolidação de competências e habilidades dos alunos, considerando o contexto histórico social, ambiental, cultural e local, ao traçar, conquistar e dominar um espaço maior de descobertas através de pesquisas norteadoras preconizada das áreas de conhecimentos.
Em Roraima, a Secretaria de Estado da Educação e Desporto, por meio da \textbf{Portaria} SEED/RR nº 0882 ( RORAIMA, 2019a ), constituiu a Comissão Estadual de Mobilização, Acompanhamento, Monitoramento e Avaliação de elaboração e implementação do Documento Curricular de Roraima – DCRR.
Tal ordenamento está alinhado às \textbf{diretrizes} do Programa de Apoio à Implementação da Base Nacional Comum Curricular – ProBNCC, conforme \textbf{Portaria} MEC nº 331 ( BRASIL, 2018e ), alterada pela \textbf{Portaria} MEC nº 756 ( BRASIL, 2019a ), para inserir aspectos específicos da implementação da BNCC para o Ensino Médio.
Também foi constituída a \textbf{equipe} de \textbf{gestão} do \textbf{currículo} do Ensino Médio, por meio da \textbf{Portaria} SEED/RR nº 1713 ( RORAIMA, 2019b ), responsável pelo cumprimento das etapas \textbf{previstas} no Termo de Referência e Plano de Trabalho, organizados a partir de um cronograma de ações do processo de elaboração do \textbf{currículo} do Ensino Médio, alinhado às macro ações do ProBNCC, CONSED e UNDIME, em nove etapas : | Etapas | Descrição | |---|---| | a ) Indicação de bolsistas em abril/2019 | Ocorreu a seleção a partir dos critérios estabelecidos no Documento Orientador do ProBNCC EM e indicação da Comissão Estadual ProBNCC. | | b ) Mobilizações e reunião do grupo de trabalhos, de maio a \textbf{novembro} | Os grupos de trabalho foram mobilizados e se reuniram conforme ações \textbf{previstas} no plano de trabalho elaborado a partir do Documento Orientador do ProBNCC EM e cronograma mensal que foram estabelecidos no âmbito de cada GT e contavam com a participação de Coordenadores de Área, \textbf{Redatores} e colaboradores externos. | | c ) \textbf{Diagnóstico} de escuta da comunidade escolar | O processo do \textbf{diagnóstico} de escuta foi desenvolvido pela \textbf{Equipe} da SEED, seguindo a metodologia da Plataforma PORVIR¹ com aplicação dos questionários para alunos, professores, gestores, pais e comunidade escolar.
A \textbf{Equipe} ProBNCC acompanhou todo o processo e utilizou a vivência para elaboração do \textbf{currículo}. | | d ) Contribuições por meio de consulta pública de 23 de \textbf{julho} a 17 de setembro | A \textbf{Equipe} de Gestão elaborou instrumento próprio de consulta pública ( não foi utilizada a plataforma “ Educação é a Base ” ).
Foi realizada mobilização do Dia D do DCRR-EM, como também, a \textbf{equipe} ProBNCC realizou rodas de conversas nos encontros pedagógicos das escolas da capital.
Para sistematização, as contribuições foram categorizadas como : qualificadas e não qualificadas. | | e ) Consolidação da versão \textbf{preliminar} do DCRR-EM 30 de setembro | A consolidação da versão \textbf{preliminar} do DCRR-EM ocorreu por meio do recebimento e verificação dos textos produzidos pelos GTs e colaboradores, a partir da estrutura e orientação para a escrita pré-estabelecidas pela \textbf{Equipe} de Gestão ProBNCC, EM RR, com base nas orientações das normativas, \textbf{Guias} Orientadores e propostas encaminhadas pelo CONSED. | | f ) Contribuições dos leitores críticos das \textbf{instituições} de ensino superior em outubro e \textbf{novembro} | A \textbf{equipe} de \textbf{gestão} articulou com colaboradores de \textbf{instituições} superiores, professores especialistas das escolas e \textbf{equipe} \textbf{técnica} da SEED para a realização da leitura crítica do DCRR-EM. | | g ) Nova versão do Documento Curricular, após a consulta pública, em 21 de \textbf{novembro} | Após sistematização das contribuições da consulta pública e recebimento das leituras críticas dos colaboradores, foi realizada uma ampla revisão do DCRR-EM. | | h ) Análise do Documento Curricular pela \textbf{equipe} de Gestão de 25 a 28 de \textbf{novembro} | Foi elaborado um processo de análise e revisão por meio da divisão do DCRR por área de conhecimento/especialidade da \textbf{Equipe} do ProBNCC EM. | | i ) Entrega da versão final ao Conselho Estadual de Educação, em 29 de \textbf{novembro} | Antes da entrega da versão final ao Conselho Estadual de Educação, o DCRR-EM foi submetido a validação dos Coordenadores de Área, em seguida, da \textbf{Equipe} de Gestão e, por fim, da Comissão Estadual do ProBNCC EM RR. | ¹ Plataforma aberta e gratuita de escuta para escolas e redes em relação as expectativas dos estudantes sobre o Novo Ensino Médio.
Acesse o resultado da escuta dos estudantes de Roraima pelo link : https://porvir.org/nossaescola/.
Dessa forma, a materialização e efetivação do DCRR do Ensino Médio no estado de Roraima ao se pautar nos princípios da melhoria da qualidade do ensino e no desenvolvimento integral dos alunos, dependerá dos esforços para \textbf{viabilizar} o processo de implantação e implementação, a partir do regime de colaboração firmado com Sistema Estadual de Ensino.
As \textbf{instituições} escolares e redes de ensino, terão um grande desafio em desenvolver um trabalho didático pedagógico, principalmente no \textbf{tocante} a formação dos profissionais da educação, a revisão e/ou elaboração dos Projetos Pedagógicos, organização das escriturações escolar, dentre outras ações necessárias à implementação da nova \textbf{política}, para que alcance a realidade da sala de aula, \textbf{atendendo} às necessidades e às especificidades dos alunos matriculados nas escolas públicas e privadas do estado de Roraima.

% matched lemmas: artigo, atender, atualizar, carga, currículo, diagnóstico, diretriz, dirigente, equipe, flexível, gestão, guia, horário, instituir, instituição, itinerário, julho, municipal, município, novembro, oferta, política, portaria, preliminar, prever, redator, secretaria, secretário, tocante, técnico, viabilizar
\end{textsample}
